\documentclass[11pt,a4paper]{article}
\usepackage[margin=1.8cm]{geometry}
\usepackage{amsmath,amssymb,tabularx,array,needspace,fontspec,enumitem,cancel,booktabs}
\newcommand{\tb}{\cancel{||||}\,}
\setmainfont{FreeSerif}
\newfontfamily\urdufont[Script=Arabic]{FreeSerif}
\TeXXeTstate=1
\newcommand{\urdu}[1]{{\urdufont\beginR #1\endR}}
\setlength{\parindent}{0pt}
\setlength{\parskip}{0pt}
\renewcommand{\arraystretch}{1.15}
\pagestyle{plain}

% \lqn{number}{content}
\newcommand{\lqn}[2]{\Needspace{6\baselineskip}\par\medskip\noindent\hangindent=2.0em\hangafter=1\makebox[2.0em][l]{\textbf{#1.}}#2\par}
\begin{document}
\begin{center}
  {\Huge\bfseries The Hope Science Academy}\\[3pt]
  {\small Excellence is our standard not goal}\\[2pt]
  {\small 247-E-1, (Opposite LDA School) Johar Town, Lahore. PH: 0300-8194789, 0322-4157001}
\end{center}

\medskip
\noindent
\begin{tabularx}{\linewidth}{|X|X|X|X|}
\hline
\textbf{Student Name:} & \textbf{Roll Num:} & \textbf{Class:} 12th & \textbf{Paper Code:} \\ \hline
\textbf{Subject Name:} Principles of Accounting & \textbf{Time Allowed:} & \textbf{Total Marks:} 595 & \textbf{Exam Date:} 02-Nov-2026 \\ \hline
\multicolumn{4}{|l|}{\textbf{Exam Syllabus:} FULL BOOK} \\ \hline
\end{tabularx}

\bigskip
\noindent\textbf{Q3. Write detailed answers of the following questions.}\hfill\textbf{1X121=121}

\lqn{1}{From the below you are required to ascertain the profit ord loss made by the trader for the year ending 31.12.2022.\\
A trader keeps his books by the single entry method. His position on 31\textsuperscript{st}\\
December 2021 was as follows:\\
Cash at Bank Rs. 6,000; stock Rs. 40,000; Debtors Rs. 60,000; machinery Rs. 100,000; and creditors Rs. 46,000. His position on 31\textsuperscript{st}\\
December 2022 was as follows:\\
Cash at Bank Rs. 84,000; stock Rs. 50,000; Debtors Rs. 90,000; machinery Rs. 90,000; and creditors Rs. 50,000.\\
During the year, the trader introduced Rs. 20,000 as further capital in business and withdrew Rs. 600 per month.\par
}
\lqn{2}{Khubaib keeps his books by single system. His capital on 30\textsuperscript{th}  June 2022 was Rs. 76,000. His capital on 1\textsuperscript{st}  July 2021 was Rs. 40,000. He withdrew for his household purpose Rs. 1,200 per month. He sold his investment worth Rs. 8,000 at 3\% premium and brought that money into the business. Calculate the profit of Khubaib for the year.\par
}
\lqn{3}{You are required to prepare the statement of profit and loss.\\
Following balances appear in Wasim’s book which are kept in the single entry basis.\\
Capital on 31\textsuperscript{st}\\
December 2022 Rs. 192,000\\
Capital on 1\textsuperscript{st}\\
January 2022 Rs. 200,000\\
He withdrew for his domestic purpose Rs. 34,800. Wasim sold his laptop Rs. 40,000 at 10\% loss and brought the amount into the business.\par
}
\lqn{4}{Ansar commenced business on 1.7.2021 with a capital of Rs. 40,000. During the year ending 30.6.20227, he borrowed Rs. 10000 from Mr. A and introduced a further capital of Rs. 42,000. Furniture worth Rs. 6,000 was purchased. He had drown from the business a sum of Rs. 6,000 during the year. On 30.6.2022 the position was as follows:\\
Cash in hand Rs. 900; Cash at bank Rs. 5,000; sundry debtors Rs. 8,1000; stock Rs. 14,000; bills receivable Rs. 3,600; sundry creditors Rs. 1,500.\\
Furniture is to be depreciated by 10 percent. Ascertain the profit or loss made by Ansar during the year ended 30.6.20227.\par
}
\lqn{5}{Nawaz commenced business on 1\textsuperscript{st}  January, 2022 with a capital of Rs. 50,000. He immediately bought furniture and fixture for Rs. 10,000. On 30\textsuperscript{th}  June, 2022, he borrowed Rs. 25,000 from his wife at 9\% per annum (interest not yet paid) and introduced a further capital of his own amounting Rs. 7,500. Nawaz drew at the rate of Rs. 1,500 per month at the end of each month for household purpose. On 31\textsuperscript{st}  December, 2022 his position was as follows:\\
Cash in hand, Rs. 1,000; Cash at bank, Rs.13,000; sundry debtors, Rs. 24,000; stock Rs. 34,000; Bills receivable Rs. 8,000; sundry creditors, Rs. 2,500; and owing for rent Rs. 750. Furniture and fixture are to be depreciated by 10\%. Ascertain the profit and loss of Nawaz during 2022.\par
}
\lqn{6}{Masood keeps his books by single entry system. He gives you the following information from which you are required to ascertain his profit or loss during 2022.\par
\par\smallskip\noindent\hspace*{1.2em}\begin{tabular}{|ccc|}
\hline
 & Jan.1, 2022 & Dec.31, 2022 \\
 & Rs. & Rs. \\
Cash at bank & 8,400 & 6,500 \\
Cash in hand & - & 200 \\
Debtors & 52,000 & 85,000 \\
Creditors & 16,000 & 20,500 \\
Stock & 16,500 & 18,000 \\
Plant & 20,500 & 21,000 \\
Furniture & 1,500 & 1,500 \\
\hline
\end{tabular}\par\smallskip
Masood had withdrawn Rs. 25,000 during the year but had introduced fresh capital of Rs. 7,000 on 1\textsuperscript{st}  July. A provision of 10\% on debtors is necessary. Write off depreciation on plants at 15\%, Interest on capital is to be allowed at 6\% P.A.\par
}
\lqn{7}{Mr. Saleem Hassan commenced business on 1\textsuperscript{st}  Jan., 2022, with a capital of Rs. 100,000. He immediately purchased furniture of Rs. 24,000. On 1\textsuperscript{st}  April, 2022, he borrowed Rs. 60,000 from his friend Mr. Nasir @ 10 \% p.a. and introduced a further capital of his own amounting Rs. 38,000. He had withdrawn Rs. 3,600 at the end of each month for family expanses. On 31\textsuperscript{st}  Dec. 2022 his position was as follows:\\
Cash in hand Rs. 2,400, cash at bank Rs.31,200, sundry debtors 57,600, stock Rs. 81,600, B/R Rs. 19,200, sundry creditors Rs. 6,000, rent due Rs. 1,800.\\
Ascertain the profit or loss made by him after considering the following adjustments:\\
1) Depreciate furniture @ 10\% p.a.\\
2) Charge interest on capital @ 10\%.\\
3) Interest is unpaid on Mr. Nasir’s loan.\\
4) Write off Rs 5,600 from Sundry debtors.\par
}
\lqn{8}{A retail trade does not keep double entry records. His business Balance Sheet at 1\textsuperscript{st}  January, 2022 was as follows:\par
\par\smallskip\noindent\hspace*{1.2em}\begin{tabular}{|clcc|}
\hline
 & Balance Sheet &  &  \\
Assets & Rs. & Liabilities & Rs. \\
Cash in hand & 1,000 & Creditors & 8,000 \\
Cash at bank & 12,000 & Capital & 24,000 \\
Stock & 10,000 &  &  \\
Debtors & 6,600 &  &  \\
Furniture & 2,400 &  &  \\
 & \underline{32,000} &  & \underline{32,000} \\
\hline
\end{tabular}\par\smallskip
His position as at the end of the year 2022 was:\\
Cash in hand Rs. 1,000; cash at bank Rs.10,000; stock Rs. 12,000; debtors Rs. 9,200; furniture Rs. 3,000 and creditors Rs. 12,000.\\
He withdrew during the year Rs. 18,000 out of which he spent Rs 12,000 on the cost of purchasing a motor car for the business.\\
Prepare necessary statements to determine the net profits for the year and a balance sheet as at 31\textsuperscript{st}  December 2022 after making the following adjustments:\\
Depreciate furniture and motor car at the rate of 10\%, write off bad debts Rs.400 and provide 5\% for the doubtful debts.\par
}
\lqn{9}{Sagheer is not writing his book properly. From the following information prepare a statement showing profit or loss and statement of affairs for the year ending 30\textsuperscript{th}  June, 2022.\par
\par\smallskip\noindent\hspace*{1.2em}\begin{tabular}{|ccc|}
\hline
 & 1.7.2021 & 30.6.2022 \\
Cash in hand & 1,800 & 5,600 \\
Debtors & 45600 & 42800 \\
Creditors & 62400 & 56800 \\
Stock & 66800 & 74800 \\
Bills receivable & 61000 & 57600 \\
Bank overdrafts & 81600 & 78400 \\
Motor van & 8400 & 8400 \\
Furniture & 6800 & 6800 \\
\hline
\end{tabular}\par\smallskip
Drawings Rs. 9600; depreciate furniture at 10\%. Write off Rs. 1600 on motor van. Provide Rs. 2000 as bad debts and 5\% as reserve on debtors. Provide reserve of Rs. 3200 on Bills receivable.\par
}
\lqn{10}{Nayyar keeps his books on single entry system. His statement of Assets and Liabilities as on 31\textsuperscript{st}  December 2022 is as follows:\par
\par\smallskip\noindent\hspace*{1.2em}\begin{tabular}{|cccc|}
\hline
Assets & Rs. & Liabilities & Rs. \\
Land and buildings & 40,000 & Sundry Creditors & 130,000 \\
Furniture \& Fixture & 6,000 & Loan from money lender & 113,200 \\
Plants \& Machinery & 110,000 & Other Liabilities & 26,800 \\
Stock & 18,400 &  &  \\
Sundry Debtors & 151,000 &  &  \\
Cash & 16,600 &  &  \\
 & \underline{342,000} &  & \underline{270,000} \\
\hline
\end{tabular}\par\smallskip
His drawings during the year amount to Rs. 6000. Loan and building are to be depreciated by 2\%, Furniture and fixtures by 5\% and plant and machinery by 10\%. Sundry debtors are to be reduced by 2\%. He has used Rs. 1600 worth of stock of his business for private purposes. During the year 2022 he sold some of his household furniture for Rs. 2,000 and paid this into the business Bank Account.  His capital at the beginning of the year was Rs. 60,000. Draw up his statement of profit or loss of the year ending 31\textsuperscript{st}  December 2022.\par
}
\lqn{11}{X and Y carrying a business in partnership keep their books by single entry. On 1\textsuperscript{st}  January 2022, the statement of their position was:\par
\par\smallskip\noindent\hspace*{1.2em}\begin{tabular}{|cccc|}
\hline
Assets & Rs. & Liabilities \& Capital & Rs. \\
Cash in hand & 270 & Bills payable & 6,460 \\
Cash at bank & 2,190 & Sundry Creditors & 20,280 \\
Bills receivable & 4,060 & Capital amounts &  \\
Sundry debtors & 48,670 & X                   73,400 &  \\
Stock & 32,850 & Y                   73,400 & 1,46,800 \\
Plants and machinery & 80,200 &  &  \\
Furniture & 5,300 &  &  \\
 & \underline{1,73,540} &  & \underline{1,73,540} \\
\hline
\end{tabular}\par\smallskip
The following was the state of affairs on 31\textsuperscript{st}  December, 2022.\\
Cash in hand, Rs. 400, cash at bank Rs. 5,810, debtors Rs. 56,280, bills receivable Rs. 6,840, stock Rs. 36,730; creditors Rs. 21,470, bills payable Rs. 5,950.\\
The partners have drawn 4,800 each and were entitled to interest on capital at 6\% per annum. No interest was payable on drawings. It was agreed to depreciate plant and machinery at 10\% and furniture at 6\%.\\
Draw up a statement of profit and loss for the year ended 31\textsuperscript{st}  December 2022, and also a statement of affairs as on that date.\par
}
\lqn{12}{Barkat keeps his books by single entry principle. On 1\textsuperscript{st}  January, 2022 his assets and liabilities were as follows:\\
Capital Rs. 72,000, Sundry Creditors Rs. 18,500, Sundry Debtors Rs. 22,000, Machinery Rs. 34,000; Furniture Rs. 5,000; Stock Rs. 14,000.\\
The following is the cash book in respect of 2022:\par
\par\smallskip\noindent\hspace*{1.2em}\begin{tabular}{|cccc|}
\hline
Receipts & Rs. & Payments & Rs. \\
Balance b/d & 15,500 & Creditors & 22,000 \\
Received from debtors & 40,000 & Wages & 6,000 \\
Sales & 12,000 & Expenses & 12,000 \\
 &  & Drawings & 5,000 \\
 &  & Balance c/d & 22,500 \\
 & \underline{67,500} &  & \underline{67,500} \\
\hline
\end{tabular}\par\smallskip
On 31\textsuperscript{st}  December 2022 the sundry creditors were Rs 20,000; sundry debtors Rs. 28,000, Machinery Rs. 34,000; Furniture Rs. 5,000, and stock Rs. 18,000.\\
Prepare Trading and Profit or Loss Account for the year ended 31\textsuperscript{st}  December, 2022 and Balance Sheet after providing for depreciation on Machinery at 10\% and Furniture at 15\% interest on Capital at 5\% and 5\% on Debtors for doubtful debts.\par
}
\lqn{13}{From the below figures, prepare final accounts after depreciation 5\% on furniture and 10\% on business premises Rs. 3,500 for doubtful debts.\\
Mahmood keeps his books by single entry and desires you to prepare Profit or Loss Account for the year ended 30\textsuperscript{th}  September, 2022 together with a Balance Sheet as on that date, upon analyzing the cash book for the year, you find the following:\par
\par\smallskip\noindent\hspace*{1.2em}\begin{tabular}{|cccc|}
\hline
 & Rs. &  & Rs. \\
Receipts from debtors & 78,000 & Bank overdraft on 1\textsuperscript{st}  Oct., 2016 & 19,500 \\
Receipts from cash sales & 48,000 & Interest on bank overdraft & 525 \\
 &  & Manager’s salary & 4,000 \\
 &  & Staff salary & 24,000 \\
 &  & General expenses & 24,275 \\
 &  & Paid to trade Creditors & 46,000 \\
 &  & Balance with cashier on 30\textsuperscript{th}  Sept. 2017 & 7,700 \\
 & \underline{1,26,000} &  & \underline{1,26,000} \\
\hline
\end{tabular}\par\smallskip
His capital on 1\textsuperscript{st}  October, 2021 was Rs. 1,03,000, stock Rs 28,500. Debtors Rs. 68,000. Furniture Rs. 4,000. Business premises Rs. 48,000 and Creditors Rs. 26,000. On 30\textsuperscript{th}  September, 2022 his stock was Rs. 32,660. Debtors Rs. 92,000. Furniture Rs 4,000 Premises Rs. 48,000, Creditors Rs. 18,500.\par
}
\lqn{14}{Prepare trading and profit and loss account for the year ended 31\textsuperscript{st} Dec., 2022 and the balance sheet as at that date.\\
Arshad keeps his books by single entry system. He gives you the following information.\par
\par\smallskip\noindent\hspace*{1.2em}\begin{tabular}{|ccc|}
\hline
 & 1-1-2022 & 31-12-2022 \\
 & Rs. & Rs. \\
Cash in hand & 500 & - \\
Cash at bank & 2500 & 7000 \\
Debtors & 7500 & 8000 \\
Stock in trade & 45000 & 37500 \\
Furniture & 5000 & 5000 \\
Creditors & 15000 & 10000 \\
\hline
\end{tabular}\par\smallskip
During the year he received Rs. 110000 from debtors and paid Rs. 75000 to creditors. He paid Rs. 21000 for wages which includes Rs. 200 per week to himself. He paid Rs. 10000 to sundry expenses. The creditors figure Rs. 75000 include Rs. 2500 for goods of his own use. He charges 10\% depreciation on furniture.\par
}
\lqn{15}{Prepare a Trading Profit and Loss Account for the year ended 31\textsuperscript{st} December, 2022 and a balance sheet as on that date.\\
Mr. Tipu keeps his books under single entry system. The position of his business as on 1\textsuperscript{st}  Jan., 2022 was as under:\\
Sundry creditors Rs. 17,000; Freehold premises, Rs. 50,000; stock Rs. 25,000; Sundry Debtors Rs. 20,000; Furniture Rs. 2,000. An abstract of Cash book is appended below:\par
\par\smallskip\noindent\hspace*{1.2em}\begin{tabular}{|cccc|}
\hline
Receipts & Rs. & Payments & Rs. \\
Sundry Debtors & 15,000 & Overdraft ( 1.1.2016) & 10,000 \\
Cash Sales & 80,000 & Expenses & 50,000 \\
 &  & Drawings & 3,000 \\
 &  & Sundry creditors & 20,000 \\
 &  & Cash in hand & 2,000 \\
 &  & Cash at bank & 10,000 \\
 & \underline{95,000} &  & \underline{95000} \\
\hline
\end{tabular}\par\smallskip
The following additional information is available:\\
Closing stock Rs. 30,000; Closing Debtors Rs. 25,000; closing creditor Rs. 12,000. No additions were made during the year to premises and Furniture account, but they are to be depreciated at 10\% and 15\% respectively. A bad debts reserve on 2\textsuperscript{1/2} \% is to be raised.\par
}
\lqn{16}{Mr. Naeem commenced a business on 1\textsuperscript{st}  January 2022 with a capital of Rs. 70,000. He purchased furniture on 1\textsuperscript{st}  April 2022 for cash Rs. 15000. On 1\textsuperscript{st}  July he purchased premises for cash Rs. 50,000.\\
From the following details obtained from his book by single entry, you are asked to prepare a Trading and Profit and Loss Account for the year ending 31\textsuperscript{st}  December, 2022 and a balance sheet as on that date providing depreciation on furniture and premises @ 10\% p.a.\par
\par\smallskip\noindent\hspace*{1.2em}\begin{tabular}{|cc|}
\hline
Sale (Inclusive of Cash Rs. 45000) & Rs. 100,000 \\
Purchases (Exclusive of cash Rs. 10,000) & Rs. 25,000 \\
Drawings & 5,000 \\
Wedges & 8,000 \\
Rent & 6,000 \\
Salaries & 28,000 \\
Bad debts written off & 2,500 \\
\hline
\end{tabular}\par\smallskip
Naeem took worth Rs. 4000 from the business and paid 1500 to his son, but omitted to record these transactions in his book. On 31\textsuperscript{st}  December, 2022 his debtors were Rs. 29,000 and creditors Rs. 14000. Stock on that date was Rs. 20,000.\par
}
\lqn{17}{From the following receipts and payments account for the year ended 31\textsuperscript{st}  December 2016, prepare the income and expenditure account.\par
\par\smallskip\noindent\hspace*{1.2em}\begin{tabular}{|lclc|}
\hline
Receipts & Rs. & Payments & Rs. \\
Balance 1-1-2016 & 2,480 & Salaries and wedges & 1,100 \\
Subscription & 3,260 & Rents & 1,000 \\
Life membership fee & 500 & General expenses & 1,100 \\
Interest & 260 & Furniture & 1,600 \\
Admission fee & 480 & Printing Expenses & 250 \\
Misc. receipts & 100 & Repairing Charges & 300 \\
Donations for building funds & 20,000 & Balance 31-12-2016 & 21,830 \\
 & \underline{27,180} &  & \underline{27,180} \\
\hline
\end{tabular}\par\smallskip
The subscription in appears amounted to Rs. 700.\\
Depreciation on furniture is at 10\%\\
Life membership fee to be credited to revenue.\par
}
\lqn{18}{The following is the Receipts and Payments Account of young man’s society for the year ending 31\textsuperscript{st}  December, 2016. You are required to prepare income and expenditure account for the year ending 31\textsuperscript{st}  December, 2016. Sport Material Valued on 31.12 .2016 at Rs. 4,000.\\
Receipts and Payment Account\par
\par\smallskip\noindent\hspace*{1.2em}\begin{tabular}{lcl}
Receipts & Rs. & Payments \\
Balance 1-1-2016 & 3,485 & Books \\
Entrance fees & 650 & Printings and Stationaries \\
Donations & 6,000 & Newspapers \\
Subscriptions & 8,565 & Sports material purchased \\
Interest on investments & 200 & Repairs \\
Sale of furniture (Book Value Rs. 1000) & 685 & Investments \\
Sale of old papers & 465 & Furniture \\
Receipts for recreations & 865 & Salaries \\
Miscellaneous Receipts & 125 & Balance 31-2-2016 \\
 & \underline{\underline{21,040}} &  \\
\end{tabular}\par\smallskip
}
\lqn{19}{You are required to prepare an Income and Expenditure Account for the year ended December 31, 2016.\\
The following particulars are obtained from the books of Young men’s Association, for the year ended 31, December 2016.\\
Receipts and payments account for the year ended December 31, 2016.\par
\par\smallskip\noindent\hspace*{1.2em}\begin{tabular}{|lclc|}
\hline
Receipts & Rs. & Payments & Rs. \\
Balance b/d & 7000 & Salaries & 3000 \\
Subscription & 9000 & Printing and Stationaries & 300 \\
Entrance fees & 1000 & Postage and telephone & 1800 \\
Life membership fees & 5000 & Municipal Taxes & 1600 \\
Annual dinner & 6000 & Electricity charges & 1200 \\
Hall rent & 3000 & Annual dinner charges & 4800 \\
Miscellaneous Income & 600 & Fixed deposit with ABL bank & 12000 \\
 &  & Balance c/d & 6900 \\
 & \underline{31,600} &  & \underline{31,600} \\
\hline
\end{tabular}\par\smallskip
The Annual Subscription of the association members is Rs. 10, and 100 members are in areas during this year.\\
Municipal taxes were paid up to Match 31, 2017.\par
}
\lqn{20}{Following are the Receipts and Payments account of an Officer’s Club for the year ended\\
31\textsuperscript{st} December 2016:\par
\par\smallskip\noindent\hspace*{1.2em}\begin{tabular}{|lclc|}
\hline
 & Rs. &  & Rs. \\
Opening Balance & 8000 & Salaries & 3600 \\
Subscription & 21500 & Telephone Charges & 750 \\
Entrance fees & 1300 & Insurance & 1200 \\
Sale of old Magazines & 50 & Miscellaneous Expenditures & 480 \\
Interest on Investments & 500 & Subscriptions to newspaper & 640 \\
Other Receipts & 800 & Purchase of Sports Material & 6000 \\
 &  & Closing Balance & 19480 \\
 & 32,150 &  & 32,150 \\
\hline
\end{tabular}\par\smallskip
2015 December 31\textsuperscript{st}  --- Balance of investments Rs. 10,000, Buildings Rs. 25,000; Subscriptions due 3,000, The Balance as on 31\textsuperscript{st}  December 2016 are --- subscriptions due Rs. 2,500. Depreciate buildings by 5\%. Interest due on investments Rs. 500. Prepare income and expenditure account.\par
}
\lqn{21}{From the following particulars prepare Income and Expenditure Account for the year ended\\
31\textsuperscript{st} December 2016.\\
The following is the receipts and payments Account of a Sports Club for the year ended\\
31\textsuperscript{st}  December, 2016.\par
\par\smallskip\noindent\hspace*{1.2em}\begin{tabular}{|lclc|}
\hline
Receipts and Payments account &  &  &  \\
Dr. &  &  & Cr. \\
Balance 1\textsuperscript{st}  Jan. 2016 & 2800 & Secretary’s Salary & 3600 \\
Entrance fees & 500 & Wages to ground men (c) & 4700 \\
Subscriptions (a) & 8700 & Ground rent & 200 \\
Proceeds to concerts & 2000 & Printing and Postage & 200 \\
 &  & Sundry repairs & 200 \\
 &  & Balance on 31-12-2016 & 5100 \\
 & \underline{14000} &  & \underline{14000} \\
\hline
\end{tabular}\par\smallskip
(a) This item includes Subscriptions outstanding brought over from previous year Rs. 400.\\
(b) This item includes Rs. 100 applicable to the previous year.\\
Entrance fees are to be capitalized. The outstanding liabilities on 31\textsuperscript{st}  December, 2016 were: Sundry repairs Rs. 50\par
}
\lqn{22}{From the particulars given below prepare the Income and Expenditure Account and the Balance Sheet of a music club started on 1-4-2016 for the year ended on 31\textsuperscript{st}  March 2017:\par
\par\smallskip\noindent\hspace*{1.2em}\begin{tabular}{lcl}
Receipts & Rs. & Payments \\
Donations & 12000 & Land and buildings \\
Entrance fees & 3000 & Concert expenses \\
Subscriptions & 6000 & Purchase of public address system \\
Rent of a part of a building let out & 600 & Salaries \\
Annual grant from the State Government & 9600 & Annual reward for the Musician of the \\
 &  & Cash in hand \\
\end{tabular}\par\smallskip
Cash at bank\\
31,200\\
One-third of the entrance fees received was to be credited to revenue and donations is to be capitalized. Depreciate land and buildings by 5\% and public address system by 2\% outstanding subscriptions for the year 2016-2017 are Rs. 200 and advance receipt of subscriptions on account of the year 2017-2018 is Rs. 400. Salaries unpaid for the year 2016-2017 amount Rs. 125.\par
}
\lqn{23}{Below is given the Receipts and Payments Account of Leo’s club. Lahore, Prepare income and expenditure account as on 31\textsuperscript{st}  December 2016 and a balance sheet as on that date, from it and subjoined information.\par
\par\smallskip\noindent\hspace*{1.2em}\begin{tabular}{|lclc|}
\hline
 & Rs. &  & Rs. \\
Balance b/d & 2476 & General expenses & 1084 \\
Annual subscription & 3260 & Salaries & 1100 \\
Life membership fees & 500 & Furniture & 1600 \\
Entrance fees & 400 & Rent & 1120 \\
Miscellaneous Receipts & 466 & Printing & 250 \\
Interest & 80 & Repairs & 300 \\
 &  & Balance c/d & 1728 \\
 & \underline{7182} &  & \underline{7182} \\
\hline
\end{tabular}\par\smallskip
On 31\textsuperscript{st}  December 2016 the annual subscriptions in areas amounted to Rs. 700. Rs. 200 is to be written off as the depreciation on furniture. Capitalize the life membership fees.\par
}
\lqn{24}{Given below is the Receipts and Payments account of Eagle club for the year ending\\
31\textsuperscript{st}  December 2016:\par
\par\smallskip\noindent\hspace*{1.2em}\begin{tabular}{|lclc|}
\hline
 & Rs. &  & Rs. \\
Balance b/d & 10,250 & Salaries & 6,000 \\
Subscriptions: &  & General Expenses & 750 \\
2015 & 400 & Drama expenses & 4,500 \\
2016 & 20,500 & Newspapers etc., & 1,500 \\
2017 & 600 & Municipal Taxes & 400 \\
Donations & 5,400 & Charity & 3,500 \\
Proceeds of drama tickets & 9,500 & Investments & 20,000 \\
Sale of waste papers & 450 & Electricity charges & 1,450 \\
 &  & Balance c/d & 9,000 \\
 & \underline{47,100} &  & \underline{47,100} \\
\hline
\end{tabular}\par\smallskip
Prepare the club’s Income and expenditure account for the year ended 31\textsuperscript{st}  December, 2016 and its Balance sheet as on that date, after taking the following information into account:\\
(a)  There are 500 members, each paying an annual subscription of Rs. 50.\\
(b) Municipal taxes amounting to Rs. 400 per annum have been paid up to 31\textsuperscript{st}  March 2017, and Rs. 500 for salaries is outstanding.\\
(c)  Buildings stood in the books at Rs. 50,000 and it is required to write off depreciation at 5\%.\\
(d)  Three Percent interest has accrued on investments for five months.\par
}
\lqn{25}{From the following prepare an income and expenditure account for the year ended 31\textsuperscript{st}  March, 2017 and a Balance Sheet for the same date:\par
\par\smallskip\noindent\hspace*{1.2em}\begin{tabular}{|lclc|}
\hline
2016 & Rs. & 2017 & Rs. \\
April 1 &  & March 31 &  \\
Balance: &  & Salaries & 3600 \\
Cash at Bank & 455 & Rent & 600 \\
Cash in Office & 55 & Printing and stationaries & 145 \\
2017 &  & Postage & 25 \\
March 31 &  & Bicycle (purchased) & 95 \\
Subscription & 3000 & Prize Bond & 680 \\
(Including Rs. 200 for 2017 -18) &  & Balance: &  \\
Interest on investments &  &  &  \\
 & 1500 & Cash in office & 12 \\
(cost of investments Rs. 30,000) &  &  &  \\
Bank interests & 10 & Cash at Bank & 113 \\
Sale of Scooter & 250 &  &  \\
 & \underline{5,270} &  & \underline{5,270} \\
\hline
\end{tabular}\par\smallskip
Subscriptions include Rs. 120 for 2015- 2016. Also rent includes Rs. 50 paid for March 2016. Subscriptions Amounting to Rs. 150 have still to be collected for the year 2016- 17. Rent for March 2017 has still to be paid and Rs. 25 are outstanding against a stationary bill. The book value of the scooter was Rs. 320.\par
}
\lqn{26}{The following was the Receipts and Payments Account of Model Town Literary Society for the year ending 31\textsuperscript{st}  December, 2016.\\
Receipts and Payments Account\par
\par\smallskip\noindent\hspace*{1.2em}\begin{tabular}{lcl}
 & Rs. &  \\
Balance b/d & 500 & Rents and taxes \\
Entrance fees & 1500 & Wages \\
Subscription & 10000 & Lightning Charges \\
Life membership fees & 500 & Lectures fees \\
Donations & 3500 & Books \\
Interest on investments & 100 & Office expenses \\
Receipts from lectures & 3000 & Fixed deposits (On 1-7-2016 @ 10\% ) \\
 &  & Furniture \\
\end{tabular}\par\smallskip
Cash in hand\\
Cash at bank\\
19,100\\
At the beginning of the year society had the following assets: Books Rs. 4000; furniture Rs. 1000; Investments, Rs. 1000.\\
Subscriptions receivable at the beginning of the year was Rs. 1000, at the end of the year is Rs, 500. Rent was outstanding for 6 months (Rs. 500) both at the beginning and at the end of the year. Capitalize entrance fees and life members’ fees. Depreciate Rs. 300 on furniture and Rs. 100 on books. Interest on fixed deposit was not yet received. You are required to prepare Incomes and Expenditure Account for the year ending 31\textsuperscript{st}  December 2016 and the balance sheet as on that date.\par
}
\lqn{27}{The following is the statement of Receipts and Payments of Mumtaz Hospital for the year ending 31\textsuperscript{st}  December 2016.\par
\par\smallskip\noindent\hspace*{1.2em}\begin{tabular}{|lclc|}
\hline
Opening Balance & Rs. & Payments & Rs. \\
Cash & 500 & Salaries & 23,000 \\
Bank & 8000 & Diet Expenses & 2,500 \\
Govt. Security & 1,80,000 & Surgery Expenses & 1,200 \\
Receipts: &  & Office Expenses & 700 \\
Subscriptions & 25,000 & Miscellaneous Expenses & 100 \\
Interest & 9,000 & Closing Balances: &  \\
Donations & 4,000 & Bank & 18,000 \\
Miscellaneous & 300 & Cash & 1,300 \\
 &  & Govt. Security & 1,80,000 \\
 & \underline{2,26,800} &  & \underline{2,26,800} \\
\hline
\end{tabular}\par\smallskip
You are asked to prepare the Income and Expenditure Account for the year and balance sheet on 1\textsuperscript{st}  January 2016. The other assets on 1\textsuperscript{st}  January 2016 were: Furniture, Rs. 200; Land, Rs. 50000; Building Rs.1,50,000; Instruments, Rs. 3500. The Govt. Securities of the face value of Rs. 2,00,000 (cost Rs. 1,80,000). The subscriptions received include Rs. 10,000 for the year 2015 but Rs. 7000 is outstanding for 2016 Interest received includes Rs. 2000 for 2015 but Rs. 2300 is outstanding for 2016.\par
}
\lqn{28}{The following is the account particulars of Golf Club for the year 2016. From these, you are required to prepare Receipts and Payments Account for the above said year:\par
\par\smallskip\noindent\hspace*{1.2em}\begin{tabular}{|lc|}
\hline
 & Rs. \\
Entrance fees Received & 2500 \\
Subscriptions: &  \\
Received during the year 2016 & 19800 \\
2015 arrears & 1750 \\
2016 arrears & 1650 \\
Received in advanced in 2015 & 550 \\
Received in advanced in 2016 & 450 \\
Green fee received during the year --- 2016 & 600 \\
Expanses Paid – 2016 & 27,250 \\
Expanses outstanding --- 2016 & 1470 \\
Expanses outstanding in 2015 & 1250 \\
Surplus in bar account & 7500 \\
Secretary’s Honorarium & 1500 \\
Cash in hand -- opening balance & 2350 \\
Locker rent received & 500 \\
Investments purchased during the year cost & 1100 \\
\hline
\end{tabular}\par\smallskip
}
\lqn{29}{From the following receipts and payments account for the year ended 31\textsuperscript{st}  December 2022, prepare the income and expenditure account.\par
\par\smallskip\noindent\hspace*{1.2em}\begin{tabular}{|cccc|}
\hline
Receipts & Rs. & Payments & Rs. \\
Balance 1-1-2022 & 4,960 & Salaries and wedges & 2,200 \\
Subscription & 6,520 & Rents & 2,000 \\
Life membership fee & 1,000 & General expenses & 2,200 \\
Interest & 760 & Furniture & 3,200 \\
Admission fee & 960 & Printing Expenses & 500 \\
Misc. receipts & 200 & Repairing Charges & 600 \\
Donations for building funds & 40,000 & Balance 31-12-2016 & 43,660 \\
 & \underline{54,360} &  & \underline{54,360} \\
\hline
\end{tabular}\par\smallskip
The subscription in appears amounted to Rs. 1,400.\\
Depreciation on furniture is at 10\%\\
Life membership fee to be credited to revenue.\par
}
\lqn{30}{Pass the journal entries in the books of Bilal and prepare the consignment account.\\
Bilal consigned computers to Nasir for Rs. 100,000. He paid sundry expenses Rs. 5000. Nasir sold the whole consignment of Rs. 1,50,000. Nasir spent Rs. 12,500 for go down rent. Nasir is entitled to receive 10\% commission on sales.\par
}
\lqn{31}{Pass the journal entries in the books of consignor and prepare necessary accounts.\\
Aqeel sent on consignment of 10 refrigerators to Sufyan costing 10,000 each, and paid transportation charges Rs. 1000. Consignee spent Rs. 2000 for storage and insurance. Sufyan sold all the refrigerators at Rs. 12,000 each. Consignee settled the account after deducting 15\% commission on sale.\par
}
\lqn{32}{Prepare consignment account in books of Imran.\\
Imran consigned 500 sims to Rameez costing Rs. 20 each and paid Rs. 250 as freight. Rameez sold 350 sims at Rs. 30 each. Rameez commission is 5\% on sales\par
}
\lqn{33}{Pass the journal entries and consignment account in the books of Karim \& Co.\\
Karim \& Co. of Karachi sent on consignment to Salman and Sons. of Lahore 60 cases of cutlery goods costing Rs 175 per case. Expenses incurred by the consignor at Karachi are:\\
Freight at Rs. 275; Insurance Rs. 55; Loading charges Rs. 20.\\
Karim \& Co. Draw on Salman \& Sons. 2 months bill at sight of Rs, 7000 which the letter accepts. The charges paid by Salman \& Sons at Lahore were: unloading charges Rs. 30; storage Rs. 85; Insurance Rs. 15, Commission is payable to Salman \& Sons. at 2\% on all sales in addition to 1\textsuperscript{1/2} \% delcredere Commission.\\
The Consignee sells all the cases @ Rs. 240 per case.\par
}
\lqn{34}{Rauf \& Co. received bank remittance for the balance due from Farooq on the consignment. Pass the Journal entries and the necessary ledger accounts in the books of consignor and consignee. Rauf and Co. of Lahore consigned to Farooq of Gujrat, 60 boxes of toys at Rs 400 per box.\\
The Consignor paid carriage and insurance of Rs. 2,000. Farooq sent 10,000 at advance payment. Rauf \& Co. received an account sale from Farooq Containing the following Particular:\\
(a)  Gross sales proceeds are Rs. 34,000.\\
(b)  Transportation and warehouses charges are Rs. 1,500.\\
(c)  Commission at 15\% on gross proceeds.\par
}
\lqn{35}{You are asked to give the necessary entries in the books of consignor to record the above transactions.\\
Global Co. of Karachi sent 1,000 cases of Homeopathic Medicines to Zahid \& Bros. of Rawalpindi at Rs. 100 per case. Expenses on the consignment incurred by the consignor amounted of Rs. 3,000. Zahid \& Bros. in Rawalpindi were working as delcredere agents. Their commission is 5\%; Delcredere commission is 7\textsuperscript{1/2}  \%.\\
Exactly after six months an ‘Account Sales’ was received by Global Co. giving the following information:\\
(a)  Sales process to 750 cases Rs. 1,50,000.\\
(b) Stock of unsold goods in hand, 250 cases\\
(c) Consignee’s expenses amounted to Rs. 1,200\\
(d) A bank draft for Rs. 1,00,000 was sent by Zahid Bros. along with the account sales.\par
}
\lqn{36}{You are required to prepare the consignment account, Zeeshan’s account and goods sent on consignment account in the books of P \& Co.\\
On 1\textsuperscript{st}  July 2016 P \& Co. sends 300 cases of medicines to Zeeshan on consignment basis. The cost of each case was 1,000. P \& Co. incurred the following expenses on the consignment: Rs. 750 on railway freight and Rs. 450 on insurance.\\
200 cases were sold by Zeeshan uniformly at Rs. 2000 per cases. The expenses of Zeeshan was amounted to Rs. 1,200. He sent an account sale on 31\textsuperscript{st}  December 2016 which revealed that he charged 5\% as ordinary Commission and 10\% as delcreder commission. He sent as bank draft for Rs. 90,000.\par
}
\lqn{37}{Show these transactions in the necessary ledgers of both the parties.\\
On 1\textsuperscript{st}  March 2016 Karim Bux consigned cloth to the value of Rs. 6,50,000 at cost to Afzal \& Sons, his agent in Multan. Who accepted a 3 months drafts for Rs. 7,80,000.\\
Karim Bux paid Rs.35,000 as freight and Rs. 6,500 as insurance on 30\textsuperscript{th}  September 2016, Afzal \& Sons notified Karim Bux that 80\% of the goods had been sold for Rs. 8,24,000. and that the selling expenses were Rs. 12,100.\\
Afzal \& Sons’ were entitled to a commission at the rate of 10\% on sales.\par
}
\lqn{38}{Mr. M in Sialkot consigned to Mr. R in Swat, 200 cases of candles at 300 per case. Rs. 60,000, paying freight charges Rs. 5,000 at draws a bill on Mr. R for Rs. 30,000. 150 cases of this consignment was sold by Mr. R for 70,000 and expenses 6,000 were incurred. Mr. R is also entitled to a commission of 6\% on sales. Mr. R remitted to Mr. M, the net amount held for this account by means of a sight draft. Show the ledger account in the books of Mr. M.\par
}
\lqn{39}{Prepare consignment account, goods sent on consignment account and Mr. D’s account in the books of Mr. C.\\
On 1\textsuperscript{st}  January 2016 Mr. C of Lahore consigned to D of Sargodha goods of sale. Mr. D is entitled to a commission of 5\%. Goods costing of Rs. 24,000 were consigned to Sargodha. The\\
expenses of consignment amount to Rs 2,000. On 31\textsuperscript{st}  March 2016 an account sales was received from Mr. D showing that he had affected sales of Rs. 24,000 in respect of 3/4 of the quantity of goods consigned to him. The expenses incurred by him amounted to Rs. 1,200. Mr. D remitted the amount due by a bank draft.\par
}
\lqn{40}{Give consignment account and Babar’s account in the books of Ahmad.\\
On 1\textsuperscript{st}  January 2016 Ahmad of Lahore consigned to Babar of Multan goods for sale. Babar is entitled to commission of 5\%. Goods sent to Babar at Rs. 144,000. The expenses of the consignment amounted Rs. 7,500. On 31\textsuperscript{st}  March 2016, an account sales was received from Babar showing that he had effected sale of Rs. 120,000 in respect of 2/3 of the quantity of goods consigned to him. His actual out of pocket expenses amounted to Rs. 8,400. Babar accepts a bill drawn by Ahmad of Rs. 40,000 and remitted the balance due from him in cash.\par
}
\lqn{41}{Show consignment account, Aslam account and goods sent on consignment account.\\
Rizwan of Gujranwala sent 100 sewing machines to Aslam of Hong Kong at Rs. 650 per machine. The consignor paid Rs. 2,000 for packing and dispatching charges. Aslam, immediately after receiving the consignment accepted a bill of Rs. 40,000. After some time, Aslam reported that 80 machines were sold for Rs. 875 each and expenses being on freight, Rs. 3,000, on godown rent, Rs. 250 and on insurance Rs. 500. Aslam is entitled to a commission of 9\% on sales. Due to insolvency of a customer who purchased 4 machines failed to pay anything.\par
}
\lqn{42}{On 1\textsuperscript{st}  Jan. 2016 Razzaq \& Co. of Lahore consigned to Mushtaq \& Co. of Mardan 100 cycles at Rs. 500 each. Razzaq \& Co. had paid Rs. 200 for packing etc., Rs. 50 for insurance and Rs. 250 for carriage. On 1\textsuperscript{st}  March 2016, Mushtaq \& Co. sold 75 cycles for Rs. 45,000, the expenses thereon being Rs. 2000. Mushtaq \& Co. are entitled for commission of 5\% and 2\% (delcredere) on sales and they remitted Rs. 30,000 on account. One customer who purchased 5 cycles failed to pay his debts due to his insolvency. Pass the necessary entries and prepare the necessary ledger accounts in books of Razaaq \& Co.\par
}
\lqn{43}{Prepare necessary ledger accounts in the books of the consignor.\\
Zaheer \& Sons of Karachi dispatched a consignment of 400 cases of Electric bulbs to Hammad of Kohat on 1\textsuperscript{st}  December 2016. The cost price was quoted at Rs. 1,80,000. On 25\textsuperscript{th}  December Hammad sent an account sales with the following details:\\
(a)  75\% of goods were sold for Rs. 135,000.\\
(b)  Hammad paid unloading charges Rs. 1200.\\
(c)  Advertisement and selling expenses paid by Hammad amounted to Rs. 700.\\
(d)  Hammad is to receive a commission of 5\% on sales.\\
At the time of dispatch Zaheer \& Sons incurred an expenditures of Rs. 2,000 towards freight and insurance charges.\par
}
\lqn{44}{Normal Loss of Goods : On 1\textsuperscript{st}  Jan. 2016, Mohsin Coal Co. of Quetta consigned to Ramzan Ltd. of Lahore 10,000 tons of coal. The cost of coal was Rs. 40 per ton and railway charges were Rs. 10 per ton. On 25\textsuperscript{th}  Feb., an account sales was received from Ramzan Ltd. showing that 5000 tons of coal sold @ Rs. 80 per ton; sales expenses Rs. 4000, Insurance Rs. 1000, commission 5\% on sales. The agent enclosed a bill for the sales proceeds less expenses and reported a shortage of 100 tons on the whole consignment. Show the consignment account in books of Mohsin coal Co.\par
}
\lqn{45}{Give the consignment account on the books of M/s Naeem motors.\\
Abnormal Loss:  M/s Naeem motors sent 100 cars to Haseeb on consignment. The cost of each car was Rs. 12,000. The expenses of M/s Naeem motors were:\\
Freight Rs.  7,000, Insurance Rs. 3000.\\
During transit one car was destroyed and the insurance company admitted Rs. 9000 towards that claim.\\
Habeeb sold 7 cars at Rs. 15,000. each and paid for storage and insurance Rs. 3,400. Haseeb then accepted a bill for Rs. 90,000 at 3 months drawn by M/s Naeem motors which they discounted immediately with their bank at 6\% p.a. It was agreed that Haseeb is to get 5\% commission.\par
}
\lqn{46}{Bilal of Lahore consigned 100 cases of candles to Abdullah of Mardan which cost him Rs. 50 per case. He incurred the following costs: packing Rs. 100; Carriage Rs. 150 and Railway freight Rs. 200. Some of the cases were damaged in transit and Abdullah took delivery of 90 cases only. Abdullah spent Rs. 50 for cartage and Rs. 300 for godown rent and sold the consignment at Rs. 60 per case. He sent the net amount to Bilal after deducting his expenses and commission at the rate of 5\% on the sale proceeds together with his account sales. Bilal also received Rs. 300 from Railway for damages. Show how the consignment account would appear in books of Bilal.\par
}
\lqn{47}{Show how the figures would appear in the balance sheet of the company.\\
Sohail Ltd. was registered with an authorized capital of Rs. 40,00,000 divided into 400,000 ordinary shares of Rs. 10 each. The company issued 200,000 ordinary shares to general public for subscription at par value of 1\textsuperscript{st}  July 2016. Subscription date was fixed as 15\textsuperscript{th}  July. Pass the entries in the books of the company if:\\
Applications for Rs. 200,000 shares were received.\\
Shares were allotted on 31\textsuperscript{st}  July 2016.\par
}
\lqn{48}{Show how the figures would appear in the balance sheet of the company.\\
Nouman Ltd. was registered with an authorized capital of Rs. 40,00,000 divided into 400,000 ordinary share of Rs. 10 each. The company issued 200,000 ordinary shares to general public for subscription as par value on 1\textsuperscript{st}  July 2016. Subscription date was fixed as 15\textsuperscript{th}  July. Pass the entries in the books of the company if:\\
Applications for Rs. 230,000 shares were received.\\
Shares were allotted on 31\textsuperscript{st}  July 2016.\par
}
\lqn{49}{Pass the journal entries in the books of the company and draft a balance sheet.\\
On June 1, 2016 Decent Ltd. issued 50,000 ordinary shares @ Rs. 10 each to the general public. On 15\textsuperscript{th}  June (the subscription date) applications were received for 45000 shares. Consequently, 5000 shares were taken up by the underwriters as per agreement. On 30\textsuperscript{th}  June, 2016 the company allotted the shares.\par
}
\lqn{50}{Record the following information in the company’s books and show the figures in the balance sheet.\\
A limited company was registered with an authorized capital of Rs. 500,000 divided into 10,000 shares of Rs. 50 each. On 1\textsuperscript{st}  November, 2016 company issued 7000 shares in the following manner:\\
1000 shares to promoters against Preliminary Expenses.\\
6000 shares to general public.\\
On 15\textsuperscript{th}  November, 2016, company received application of Rs. 6000 shares.\par
}
\lqn{51}{Pass the journal entries in the books of the company and show the balance sheet.\\
A trading corporation was formed with an authorized capital of 50,000 ordinary shares of Rs. 10 each. On 1\textsuperscript{st}  Jan., 2016 the company issued 30,000 shares to the general public and 5000 shares were issued to a vendor, Mr. Asghar from whom company had acquired machinery worth Rs. 50,000. Applications were received on 10\textsuperscript{th}  Jan. 2016 only for 25000 shares. On 20\textsuperscript{th}  Jan. 2016, shares were allotted to the applicants.\par
}
\lqn{52}{Record the following transactions in the books of company and draft the balance sheet.\\
A company issued 20,000 shares of Rs. 10 each of general public at 10\% premium. Applications were received for 35000 shares. Shares were allotted and money was refunded to the applicants of 15000 shares.\par
}
\lqn{53}{Make the necessary entries and prepare the balance sheet.\\
A company passed a resolution in a general meeting to issue 20,000 shares of Rs. 10 each at a discount of Rs. 1. Applications were received for 15000 shares. Shares were allotted\par
}
\lqn{54}{Imran textile Ltd. acquired the business of M/s Noor \& Sons. The assets and liabilities of M/s Noor \& Sons at bank values are given below:\par
\par\smallskip\noindent\hspace*{1.2em}\begin{tabular}{|lccc|}
\hline
Land And Building & 50,000 & Sundry Creditors & 10,000 \\
Machinery & 40,000 &  &  \\
Furniture & 14,000 &  &  \\
Debtors & 5000 &  &  \\
\hline
\end{tabular}\par\smallskip
The purchase consideration is to be paid by the company in fully paid up shares of Rs. 10 each.\\
Pass the journal entries if the shares are issued: (a) At Par            (b) At 10\% Diiscount             \\
 (c) At 10\% Premium\par
}
\lqn{55}{A Ltd. Company issued 1000 12\% debentures of Rs. 100 each. What entries will be made if:\\
(a) Debentures are issued as per value.\\
(b)  Debentures are issued by 5\% premium.\\
(c)  Debentures are issued by 10\% discount.\par
}
\lqn{56}{Pass journal entries, if any. Also show how, the debentures and Bank Loan will appear in the Company’s Balance Sheet.\\
Z Company Ltd. had 12\% debentures of 600,000 outstanding on 1\textsuperscript{st}  January, 2016. On 1\textsuperscript{st}  July company obtained a loan of Rs. 1,75,000 from MCB for which the company placed with bank, debentures of Rs. 200,000 as collateral securities.\par
}
\lqn{57}{Pass the journal entries for issuance of debentures:\\
(a) 6,000 debentures of Rs. 100 each issued at par of cash.\\
(b) 5,000 debentures of Rs. 100 each at discount of Rs.5 for cash.\\
(c) 4,000 debentures of Rs. 100 each at 10\% premium for cash.\\
(d) To a vendor from whom furniture purchased for Rs. 50,000; 5500 debentures of Rs. 10 each issued for settlement.\\
(e) Issuance of Rs. 180,000 debentures as collateral security for a loan of Rs. 150,000 from bank.\par
}
\lqn{58}{Ali industries Ltd. issued 5000 10\% debentures of Rs. 1000 each. What entries will be passed if:\\
1.  Debentures are issued at par and redeemable at par.\\
2.  Debentures are issued at 7\% and redeemable at par.\\
3.  Debentures are issued at 8\% premium and repayable at par.\\
4.  Debentures are issued at par and redeemable at 4\% premium.\\
5.  Debentures are issued at 5\% discount and redeemable at 4\% premium.\par
}
\lqn{59}{Pass the following entries of debentures in each case:\\
(i)  Debentures issued at Rs. 90 and redeemable at Rs. 100.\\
(ii)  Debentures issued at Rs. 100 and redeemable at Rs. 110.\\
(iii)  Debentures issued at Rs. 110 and redeemable at Rs. 100.\\
(iv)  Debentures issued at Rs. 90 and redeemable at Rs. 110.\\
(v)  Debentures issued at Rs. 100 and redeemable at Rs. 100.\par
}
\lqn{60}{Haider Ltd. acquired the assets and liabilities of Mr. Rashid valued at Rs. 500,000 and Rs. 40,000 respectively. It was agreed that the purchase consideration would be settled by issue of 10\% debentures of Rs. 100 each by Haider Ltd. pass the journal entries if the purchase consideration was:  (a)  Rs. 460,000      (b)  Rs. 480,000      (c)  Rs. 430,000\par
}
\lqn{61}{X Company Ltd. purchased the business of Naeem Bros. for a purchase consideration of Rs. 7,70,000. The book value of the assets were Rs. 782,000 and those of liabilities Rs. 30,000. The company issued debentures of Rs. 100 each to Naeem Bros. pass the journal entries in the books of the company if:\\
(a)  Debentures were issued by par.\\
(b)  Debentures were issued 10\% discount.\\
(c)  Debentures were issued 10\% premium.\par
}
\lqn{62}{Prepare the building account for the first five years.\\
A firm purchased a building costing Rs. 900,000 and it is decided to provide depreciation at 5\% per annum on fixed installation method.\par
}
\lqn{63}{Show the machinery account for first four years and find out annual depreciation.\\
A firm purchased machinery for Rs. 50,000. On 1\textsuperscript{st}  Jan. 2013. Depreciation is to provided annually according to fixed installation method. The useful life of the asset is 10 years and residual value is Rs. 10,000.\par
}
\lqn{64}{X \& Co. purchased a machinery of Rs. 70,000 on 1\textsuperscript{st}  July, 2013. They spent Rs. 8000 on its installation. Prepare the machinery account for the first four years under straight line method of depreciation. Depreciation is written off at 10\% per annum. Assume the accounts are closed every year on 31\textsuperscript{st}  December.\par
}
\lqn{65}{On 1\textsuperscript{st}  Jan., 2013 a firm purchased machinery worth Rs. 50,000. On 1\textsuperscript{st}  July, 2015 it buys additional machinery worth Rs. 10,000 and spends Rs. 1000 on erection. The accounts are closed each year on 31\textsuperscript{st}  December. Assuming the normal depreciation to the 10\% p.a., show the machinery accounts for four years under fixed installation method.\par
}
\lqn{66}{A company purchased a machinery for Rs. 100,000 and depreciation rate is 10\% p.a. show the machinery account for the first five years under the reducing balance method\par
}
\lqn{67}{A firm purchased a machine for Rs. 210,000 on 1\textsuperscript{st}  July, 2013 and spent Rs. 24,000 as wages and installation charges. Prepare the machine account for the first four years under diminishing balance method of depreciation assuming that the accounting year of the firm ends on\\
31\textsuperscript{st}  December every year. Rate of deprecation is 10\% per annum.\par
}
\lqn{68}{On 1\textsuperscript{st}  January, 2013 a firm purchased machinery of Rs. 50,000. On 1\textsuperscript{st}  July, 2015 it buys additional machinery worth Rs. 10,000 and spends Rs. 1000 on its erection. The accounts are closed each year on 31\textsuperscript{st}  December. Assuming the normal depreciation to be 10\% p.a., show the machinery account for four years under reducing installation method.\par
}
\lqn{69}{A firm purchased a second-hand truck of Rs. 50,000 on 1\textsuperscript{st}  January, 2013 and spent Rs. 20,000 on its overhauling. Depreciation is written off 10\% p.a. on the reducing balance. On 30\textsuperscript{th}  June, 2016 the truck was sold for Rs. 30,000 being unsuitable. Prepare the truck account from 2013 to 2016 assuming that accounts are closed on 31\textsuperscript{st}  December every year.\par
}
\lqn{70}{Prepare the machinery account up to 2016 in the books of company.\\
A \& Co. purchased a machinery for Rs. 160,000 on 1\textsuperscript{st}  July 2013. The books are closed on\\
31\textsuperscript{st}  December every year. On 30\textsuperscript{th}  June 2016, it was sold for Rs. 70,000 and new machinery was purchased for Rs. 180,000 on the same date. Depreciation is charged at the rate of 15\% p.a. on original cost method.\par
}
\lqn{71}{Show machinery account under written down value methods and for three years from the information given below:\\
(a) Machinery Cost Rs. 200,000.\\
(b) Estimated Life 3 years\\
(c) Rate of depreciation 33\textsuperscript{1/3} \% p.a.\\
(d) Scrape value Rs. 20,000.\par
}
\lqn{72}{Show machinery account under original cost method for three years from the information given below:\\
(a) Machinery cost Rs. 200,000\\
(b) Estimate Life 3 years\\
(c) Residual value Rs. 20,000\par
}
\lqn{73}{\\
On 1\textsuperscript{st}  July 2013, Basharat purchased machinery of Rs. 60,000. Depreciation is to be provided for at 10\% on straight line method each year. On 31\textsuperscript{st}  October, 2015 Machinery was sold for Rs. 24,000 as they became useless. On the same date he purchased a new machinery for Rs. 20,000. Prepare machinery account from 2013 to 2016. Accounts are closed on 31\textsuperscript{st}  December every year.\par
}
\lqn{74}{Show the truck account for four years.\\
On 1\textsuperscript{st}  January 2013 a firm purchased a truck for Rs. 15,000. Depreciation is charged at 20\% p.a. on the written down value method. On 31\textsuperscript{st}  December 2016 it was arranged to replace the old truck by a new one cost Rs. 18,000 and an allowance of Rs. 4500 being made from the purchase price of new one.\par
}
\lqn{75}{A transport company purchases 10 motor trucks at Rs. 90,000 each on 1\textsuperscript{st}  April, 2013. On\\
1\textsuperscript{st}  October, 2015 one of the trucks got an accident and was completely destroyed. Rs. 54,000 are received from the insurer in full settlement. On the same day another truck was purchased for the sum of Rs. 100,000. The company wrote off depreciation at 20\% on the original cost per annum and observed the calendar year as its financial year. Give the motor trucks account from 2013 to 2015.\par
}
\lqn{76}{A, B and C are partners with capitals of Rs. 20,000, Rs. 10,000 and Rs. 6,000 respectively and shares profits and losses equally. The net profit for the year 2016 before charging interest on capitals amounted Rs. 10,800. Show the amount of each partner gain from the firm (a) if no interest calculated on capital and (b) where 5\% interest on capital is brought into account, before the adjustments of profits.\par
}
\lqn{77}{You are required to prepare the profit \& loss appropriation account necessary journal entries relating to profit and loss appropriation accounts.\\
X and Y set up partnership firm on 1.1.2016. They contributed Rs. 1,50,000 and Rs. 1,20,000 respectively as their capitals and decided to share profit and loss in the ratio of 3:2. The partnership deed provided that A is to be paid a salary of Rs. 3000 per month and B a commission of Rs. 15000. It also provided that interest on capital be allowed at 6\% per annum. The drawings for the year were A Rs. 18000 and B’s Rs. 12000. Interest on drawings was charged Rs. 810 on A’s drawings and Rs. 540 on B’s drawings. The net amount of profit as per profit and loss account for year 2016 was Rs. 106980.\par
}
\lqn{78}{You are required to prepare partners accounts when capitals are fluctuating.\\
A and B are partners sharing profits and losses in the proportions of 3:2 with capital of Rs.\\
10,000 each. Each partner is entitled to 5\% interest on his capital. B is entitled to a salary Rs. 150 per month. During the year 2016, the drawings of partners in anticipation of their shares of profits and salary A Rs. 1,000 and B Rs. 1,000. The profits of the year prior to calculation of interest on capital but after charging salary of partner amounted to Rs. 8,000. The above figure of profit is before charging depreciation at 7\textsuperscript{1/2} per cent on furniture valued at Rs. 5,000 and writing off a bad debt of Rs. 150.\par
}
\lqn{79}{You are requested to prepare partners accounts when capitals are fixed.\\
A and B are partners sharing profits and losses in the proportions of 3:2 with capitals of Rs. 10,000 each. Each partner is entitled to 5\% interest on his capital. B is entitled to a salary Rs. 150 per month. During the year 2016, the drawings of partners in anticipation of their shares of profit and salary A Rs. 1,000 and B Rs. 1,000. The profits for the year prior to calculation of interest on capital but after charging salary of the partner amounted to Rs. 8,000. The above figure of profit is before charging depreciation at 7\textsuperscript{1/2} per cent on furniture valued at Rs. 5,000 and writing off a bad debt of Rs. 150.\par
}
\lqn{80}{Calculate the interest payable to A and B for the year ending 31.12.2016.\\
A and B started business on 1.1.2016 with capitals of Rs. 60,000 and 40,000 respectively. During the year A introduced Rs. 10,000 as additional capital on 1.7.2016. They withdrew Rs. 500 p.m. for household expenses in lieu of profits. Interest on capital is to be allowed @ 10\%.\par
}
\lqn{81}{A and B are partners from 1\textsuperscript{st}  January, 2016 without any partnership agreement and they introduced capital of Rs. 35,000 and Rs. 20,000 respectively.  On 1\textsuperscript{st}  July, 2016 A advances Rs. 7500 by way of loan to the firm without any agreement as to interest. The profit and loss\\
accounts for the year 2016 discloses a profit of Rs. 8,225; but the partners cannot agree upon question of interest or upon the basis of division of profits. You are providing to divide the profits between them giving reasons for your methods.\par
}
\lqn{82}{X and Y are partners sharing profits and losses in the ratio of 3:1. Z is admitted into the partnership with a 1/4\textsuperscript{th}  share in the future profits which he takes entirely from A. Z brings in Rs. 40,000 on account of capital and Rs. 24,000 on account of goodwill. Pass journal entries.\par
}
\lqn{83}{P \& Q share profits and losses in the ratio 5:3. They admit R into partnership. R brings Rs. 100,000 as capital and Rs. 30,000 as goodwill. The new profit sharing is 7:4:3. Make journal entries.\par
}
\lqn{84}{A and B are partners sharing profits in the ratio of 3:2. C is admitted as a partner and is to be given 1/6 of the profit. His share of goodwill is Rs. 40,000. Give journal entries if old partners do not withdraw the amount of goodwill and when they withdraw the amount of goodwill.\par
}
\lqn{85}{A and B are partners sharing profits in the ratio of 3:2. They admit C into the partnership. C pays Rs. 5000 as capital and Rs. 1000 as goodwill for 1/4\textsuperscript{th}  share. Pass the journal entries and calculate the value of goodwill of whole business. New profit sharing is 3:2:2.\par
}
\lqn{86}{A and B are partners sharing profits and losses in the ratio of 3:2. They admit C as a partner for 1/4\textsuperscript{th}  share in profits. C paid 25,000 as capital. Goodwill is valued at two years purchase of three years profits of Rs. 5000, Rs. 8000 and Rs. 8000.\\
Pass journal entries:\\
(1) If C pays goodwill in cash and retained.\\
(2) If C pays goodwill in cash but withdrawn.\\
(3) If goodwill is raised.\\
(4) If goodwill is raised and written off.\par
}
\lqn{87}{Salman and Rehman carrying on business in partnership and sharing profits and losses in the ratio of 3:2 and their balance sheet stood as follows:\par
\par\smallskip\noindent\hspace*{1.2em}\begin{tabular}{|cclc|}
\hline
Assets & Rs. & Liabilities & Rs. \\
Cash & 1500 & Creditors & 11800 \\
Stock & 28000 & Salman’s capital & 51450 \\
Debtors & 19500 & Rehman’s capital & 36750 \\
Furniture & 2500 &  &  \\
Machinery & 48500 &  &  \\
 & \underline{100,000} &  & \underline{100,00} \\
\hline
\end{tabular}\par\smallskip
They admit Imran into partnership and give him 1/8\textsuperscript{th}  share in future profits on the following terms:\\
(1) Goodwill of the entire firm be valued twice the average of the last three year’s profits which amounted to Rs. 21,000, Rs. 24,000 and Rs. 25,560.\\
(2) Imran to bring in cash for the amount of his share of goodwill.\\
(3) He is to bring cash Rs. 15,000 as his capital.\\
Give journal entries and draw out the balance sheet of the new firm and state the future profit sharing ratio.\par
}
\lqn{88}{Saleem and Kaleem trading in partnership and sharing profits and losses in the proportion of 2 and 1 respectively agree to take Asif into the partnership on the following terms:\\
(1) Asif should be given 1/4\textsuperscript{th}  share and he should bring Rs. 10,000 as goodwill and Rs. 128,000 as capital.\\
(2) A reserve for bad and doubtful debts should be created at 5\%.\\
(3) The value of land \& building should be brought up by Rs. 620,000.\\
(4) Stock should be taken at Rs. 322,600.\\
The following is the Balance Sheet of the firm of Saleem, Kaleem and Rahim on the date of Asif’s admission.\par
\par\smallskip\noindent\hspace*{1.2em}\begin{tabular}{|llll|}
\hline
\underline{Balance Sheet as on 31.12.2016} &  &  &  \\
Assets & Rs. & Liabilities & Rs. \\
Cash in hand & 8000 & Sundry creditors & 50,000 \\
Debtors & 252,000 & Partner’s capitals: &  \\
Stock & 360,000 & Saleem   650,000 &  \\
Land \& Buildings & 480,000 & Kaleem   400,000 & 10,50,000 \\
 & \underline{11,00,000} &  & \underline{11,00,000} \\
\hline
\end{tabular}\par\smallskip
Pass journal entries in the books of the new firm, and the revolution accoun\par
}
\lqn{89}{Rameez and Raja are partners sharing profits in the ratio of 3:2. Their balance sheet stood as under on 31-12-2016:\par
\par\smallskip\noindent\hspace*{1.2em}\begin{tabular}{|lclc|}
\hline
\underline{Balance Sheet} &  &  &  \\
Assets & Rs. & Liabilities & Rs. \\
Cash & 2000 & Creditors & 38500 \\
Stock & 15000 & Outstanding liabilities & 4000 \\
Prepaid insurance & 1500 & Capitals: &  \\
Debtors & 9000 & Rameez    29000 &  \\
Machinery & 19000 & Raja          15000 & 44,000 \\
Buildings & 35000 &  &  \\
Furniture & 5000 &  &  \\
 & \underline{86,500} &  & \underline{86,500} \\
\hline
\end{tabular}\par\smallskip
Rahim is admitted as a new partner introducing a capital of Rs. 16,000. The new profit sharing ratio is decided as 5:3:2. Rahim is unable to bring any cash for goodwill. So it is decided to raise goodwill account, amount being calculated on the basis of Rahim’s share in the profits and the capital contributed by him. Following revaluations are made:\\
(1) Stock to be depreciated by 5\%.\\
(2) Provision for doubtful debts Rs. 100.\\
(3) Furniture to be depreciated 10\%.\\
(4) Building is valued at Rs. 40,000.\\
Show the revaluation account and calculate the value of goodwill of whole business.\par
}
\lqn{90}{The balance sheet of Mr. Shabir and Mr. Majid as on 31\textsuperscript{st}  December 2016 is set out below. They share profits and losses in the ratio of 2:1.\par
\par\smallskip\noindent\hspace*{1.2em}\begin{tabular}{|llll|}
\hline
Assets & Rs. & Liabilities & Rs. \\
Freehold property & 26000 & Shabir’s capital & 40000 \\
Stock & 12000 & Majid’s capital & 30000 \\
Debtors & 60000 & General reserved & 24000 \\
Cash & 12000 & Creditors & 16000 \\
 & \underline{1,10,000} &  & \underline{1,10,000} \\
\hline
\end{tabular}\par\smallskip
They agree to admit Pasha into the firm subject to the following term and conditions:\\
(a) Pasha will bring in Rs. 21,000 of which 9,000 will be treated as his share of Goodwill to be retained in the business.\\
(b) He will be entitled 1/4\textsuperscript{th}  share of the profits of the firm.\\
(c) Fifty per cent of General Reserved is to remain as a Reserve for bad and doubtful debts.\\
Show the journal entries giving effect to the above said arrangements (including cash transactions) and prepare the opening balance sheet of the new partnership.\par
}
\lqn{91}{A and B are partners in a firm sharing profits and losses as 5:3. The position of the firm as on 31\textsuperscript{st} March 2016 is as follows:\par
\par\smallskip\noindent\hspace*{1.2em}\begin{tabular}{|lcll|}
\hline
Assets & Rs. & Liabilities & Rs. \\
Plants \& Machinery & 40,000 & Capital accounts: &  \\
Stock & 30,000 & A               30,000 &  \\
Sundry debtors & 20,000 & B               20,000 & 50,000 \\
Bills receivable & 10,000 & Sundry creditors & 15,000 \\
Cash at bank & 7500 & Bank overdrafts & 42,000 \\
 & \underline{1,07,500} &  & \underline{1,07,500} \\
\hline
\end{tabular}\par\smallskip
C now joins them on condition that he will share 3/4\textsuperscript{th}  of the future profits, the balance of profits being shared by A and B as 5:3. He introduced Rs. 40,000 by way of capital in cash and pays off the overdraft. He also pays Rs. 4,000 by way of premium for goodwill of the business and this amount is to remain in business. The partners agree to depreciate plant by 10\% and raise a reserve against sundry debtors by 5\%.\\
You are asked to journalise the entries in the books of the firm and the resultant balance sheet and also show how will the partners share future profits.\par
}
\lqn{92}{Balance Sheet as on 31-12-2016 of Akif and Asim who were sharing profits and losses in the ratio of 2:3.\par
\par\smallskip\noindent\hspace*{1.2em}\begin{tabular}{|clll|}
\hline
 & Rs. &  & Rs. \\
Cash & 42,000 & Capitals: &  \\
Furniture & 28,000 & Akif & 62,000 \\
Stock & 44,000 & Asim & 70,000 \\
Debtors & 42,000 & Creditors & 96,000 \\
Machinery & 66,000 & Profits and Lose accounts & 14,000 \\
Goodwill & 20,000 &  &  \\
 & \underline{2,42,000} &  & \underline{2,42,000} \\
\hline
\end{tabular}\par\smallskip
They admit Ahsan into partnership on the following terms:\\
(1) Furniture and machinery to be reduced by 15\%.\\
(2) The value of stock to be given at Rs. 48,000.\\
(3) Goodwill to be valued at Rs. 26,000 and Ahsan introduced Rs. 32,000 towards capital.\\
(4) Outstanding rent amounted to Rs. 1,800.\\
Prepare revaluation account, capital account of old partner, cash account and balance sheet of the new firm.\par
}
\lqn{93}{ufiqar and Fayyaz are partners in a firm sharing profits and losses as Zulfiqar 3/4 and Fayyaz 1/4 on 1\textsuperscript{st} January, 2016, their positions was as given below:\par
\par\smallskip\noindent\hspace*{1.2em}\begin{tabular}{|llll|}
\hline
Assets & Rs. & Liabilities & Rs. \\
Plant & 40,000 & Capital accounts: &  \\
Stock & 10,000 & Zulfiqar      50,000 &  \\
Debtors & 30,000 & Fayyaz      30,000 & 80,000 \\
Cash at bank & 20,000 & Sundry creditors & 20,000 \\
 & \underline{100,000} &  & \underline{100,000} \\
\hline
\end{tabular}\par\smallskip
Riaz is now to join the partnership. He agrees to pay the partners Rs. 20,000 by way of goodwill and introduce 3/5 of the combined capital of the two existing partners after depreciating plant and stock at 20\% and 10\% respectively and raising a reserve of 10\% against Sundry Debtors. The new partner is to be allowed 1/4\textsuperscript{th}  share of the profits of the firm.\\
You are asked to record the above transactions in the books of the firm and give the resultant balance sheet of the new firm.\par
}
\lqn{94}{Roshan and Shahzad are partners in a firm sharing profits in the ratio of 3:2 on 1\textsuperscript{st} January 2016, the position of the business was as follows:\par
\par\smallskip\noindent\hspace*{1.2em}\begin{tabular}{|ccll|}
\hline
Assets & Rs. & Liabilities & Rs. \\
Goodwill & 5000 & Sundry creditors & 15,000 \\
Stock & 20,000 & Capital accounts: &  \\
Plant & 25,000 & Roshan     30,000 &  \\
Debtors & 18,000 & Shehzad   25,000 & 55,000 \\
Cash & 2000 &  &  \\
 & \underline{70,000} &  & \underline{70,000} \\
\hline
\end{tabular}\par\smallskip
Shabir agrees to join the business on the following conditions:\\
(1) He will introduce Rs. 20,000 as his capital and pay Rs. 10,000 to the partners as premium for goodwill for 1/4\textsuperscript{th}  share of the future profit of the firm.\\
(2) Revaluation of the assets of the firm will be made by reducing plant account to Rs. 20,000 and stock by 10\% and by raising a provision by bad debts by 5\% of debtors.\\
You are asked to give the necessary entries in the books of account, recording the above transactions and give the balance sheet of the new firm. Goodwill is to appear at its old figure in the new balance sheet.\par
}
\lqn{95}{Mukhtar and Munir were partners in a firm sharing profits equally. Their business position as on 31\textsuperscript{st}  June 2016 was as follows:\par
\par\smallskip\noindent\hspace*{1.2em}\begin{tabular}{|lclc|}
\hline
Balance Sheet &  &  &  \\
Assets & Rs. & Liabilities & Rs. \\
Cash in hand & 150 & Sundry creditors & 6000 \\
Stock & 3600 & Bank overdraft & 1500 \\
Sundry debtors 6200 Mukhtar’s capital 2100 &  &  &  \\
Furniture & 600 & Munir’s capital & 1600 \\
Investments & 650 &  &  \\
 & \underline{11,200} &  & \underline{11,200} \\
\hline
\end{tabular}\par\smallskip
It is agreed to take Mobeen into partnership and to make the following adjustments:\\
(1) Bad debts to be written off for Rs. 1,600.\\
(2) Value of the furniture to be reduced to Rs. 400.\\
(3) Depreciate stock at 10\%.\\
(4) Write of 20\% on investments.\\
(5) Raise goodwill for Rs. 1,000.\\
Mobeen introduced Rs. 1,000 as capital for his 1/3 share. Other partners’ capitals should be adjusted according to the new partner’s capital.\\
Pass necessary journal entries and prepare the balance sheet of the new firm.\par
}
\lqn{96}{Retirement:\\
Calculation of New and Gaining Ratio:\\
A, B and C were sharing profits as ½ to A, and 3/10 to C. C retires. Calculate new profit sharing ratio of A and B\par
}
\lqn{97}{Calculate new profit sharing ratio:\\
A, B, C and D were partners sharing profits in the ratio of 5:4:3:2. A and C retires from the firm\par
}
\lqn{98}{Calculate new profit sharing ratio:\\
A, B and C are partners in a firm sharing profits \& losses in the ratio of 5:3:2. B retires. His shares were taken up by A and C in the ratio of 2:1.\par
}
\lqn{99}{Calculate new and gaining ratio:\\
A, B and C were partners sharing profits in the ratio of 3:2:1.\\
(a) When A retires   (b) When B retires          (c) When C retires.\par
}
\lqn{100}{A and B are equal partners in a firm. Partnership deed provides that any partner may retire from the firm, on the following terms, after giving 3 months’ notice in writing:\\
Retiring partner shall be paid:\\
(1) Balance of his capital and current account.\\
(2) Share of profit to the date of retirement calculated on the basis of the average profit of the three proceeding years.\\
(3) Share of the goodwill of the firm calculated twice the average of profit of the 3 proceeding completed years\\
B gave a notice on June 30, 2016 to retired and retired on September 30\textsuperscript{,}  2016 when his capital account showed a credited balance of Rs. 20,000and current account debit balance of Rs. 3000. Profit for the 3 proceeding years were Rs. 28,000, Rs. 22,000 and Rs. 16,000. Firm’s year ends on March 31\textsuperscript{st}  every year.\\
Calculate the amount due to B.\par
}
\lqn{101}{A, B and C were carrying on business in partnership sharing profits and losses in ratio of\\
3:2:1on 31\textsuperscript{st} December 2016, balance sheet of the firm stood as follows:\par
\par\smallskip\noindent\hspace*{1.2em}\begin{tabular}{|lcll|}
\hline
Balance Sheet &  &  &  \\
Assets & Rs. & Liabilities & Rs. \\
Cash & 5900 & Sundry creditors & 13,590 \\
Debtors & 8000 & Capital accounts: &  \\
Stock & 11,690 & A              15,000 &  \\
Buildings & 23,000 & B              10,000 &  \\
 &  & C              10,000 & 35,000 \\
 & \underline{48,590} &  & \underline{48,590} \\
\hline
\end{tabular}\par\smallskip
B retired on the above mentioned date on the following terms:\\
(1) Buildings be appreciated by Rs. 7000.\\
(2) Provision for bad debts be made @ 5\% on debtors.\\
(3) Goodwill of the firm be valued at Rs. 9,000 and adjustments in this respect be made without raising Goodwill Accounts.\\
(4) Rs. 5,000 be paid to B immediately and the balance due to him be treated as a loan carrying interest @ 6\% per annum.\\
Pass journal entries to record the above transactions and show the B’s capital account on retirement.\par
}
\lqn{102}{A, B and C are partners, sharing profits equally. The balance sheet on December 31\textsuperscript{st} 2016 is as follows:\\
Balance Sheet\par
\par\smallskip\noindent\hspace*{1.2em}\begin{tabular}{lclc}
Assets & Rs. & Liabilities & Rs \\
Cash at bank & 4000 & Sundry creditors & 40 \\
Bills receivable & 3000 & Capitals: &  \\
Sundry debtors 20000 &  & A & 12 \\
Less: Reserve for bad debts  1,000 & 19,000 & B & 18 \\
Stock & 18,000 & C & 75 \\
\underline{Fixtures} & 3500 & Reserve & 60 \\
 & \underline{\underline{47,500}} &  & \underline{\underline{47}} \\
\end{tabular}\par\smallskip
B retires on that date and following adjustments are to be made for the purpose:\\
(a) Goodwill of the firm is valued at Rs 12,000\\
(b) Fixtures to be Depreciated by 5\%\\
(c) Stock to be appreciated to be 10\%\\
(d) Reserve for bad debts to be increased by rs. 500.\\
Draw up the revaluation account, Capital Accounts of partners and opening balance sheet of the continuing partners.\par
}
\lqn{103}{Amjad, Anwar and Arif were partners sharing profits and losses in the ratio of 3:2:1. Their balance sheet on 31\textsuperscript{st}  December 2016 was as follows:\par
\par\smallskip\noindent\hspace*{1.2em}\begin{tabular}{|lclc|}
\hline
Balance Sheet &  &  &  \\
Assets & Rs. & Liabilities & Rs. \\
Cash at bank & 3500 & Sundry creditors & 15000 \\
Bills receivable & 2000 & Bills payable & 5000 \\
Investments & 6500 & Capitals: &  \\
Debtors & 6000 & Amjad          10,000 &  \\
Stock & 5000 & Anwar          10,000 &  \\
Furniture & 2000 & Arif               10,000 & 30,000 \\
Buildings & 25000 &  &  \\
 & \underline{50,000} &  & \underline{50,000} \\
\hline
\end{tabular}\par\smallskip
Arif retired on the above date and the partners agreed that:\\
(a) Goodwill should be calculated on the basis of two years’ profits of average of the preceding years’ profits. The profits for the year 2013, 2014 and 2015 were Rs. 16,000. Rs. 12,000and Rs. 14,000 respectively.\\
(b) Rs. 510 to be provided for doubtful debts.\\
(c) Stock to be reduced to 10\%.\\
(d) There was an appreciation in the value of building by 5\%.\\
Show the revaluation account and the balance sheet of the continuing partner.\par
}
\lqn{104}{X and Y are in partnership sharing profits and losses in the ratio of 3:2. They insure their lives jointly for Rs. 75,000 at an annual premium of Rs. 3500 to be debited to the business. Y dies 3 months after the date of last balance sheet. According to the partners deed, Y’s legal representatives are entitled to the following payments:\\
(a) His capital as per the last balance sheet.\\
(b) Interest on above capital at 6\% per annum to the date of death.\\
(c) His share of profit to the date of death calculated on the basis of the last year profit.\\
(d) His share of insurance money.\\
His drawings are to bear interest at an annual rate of 2\%on the amount irrespective of period. Y's capital as per the last balance sheet was Rs. 40,000 and his drawings to date of death were Rs. 5000. The last year profit was Rs. 30,000.\\
Draw Y’s account to be rendered to his legal representations.\par
}
\lqn{105}{A, B and C are partners sharing profits and losses in proportions of 3:2:1and their balance sheet on 31\textsuperscript{st} December, 2016 was as follows:\par
\par\smallskip\noindent\hspace*{1.2em}\begin{tabular}{|lclc|}
\hline
Balance Sheet &  &  &  \\
Assets & Rs. & Liabilities & Rs. \\
Cash in hand & 250 & Bill payable & 7560 \\
Cash at bank & 960 & Creditors & 12300 \\
Bills receivable & 3,300 & General reserves & 3000 \\
Debtors & 7,450 & Capitals: &  \\
Stock & 12,470 & A               10000 &  \\
Investments & 10,430 & B               6000 &  \\
Buildings & 8,000 & C               4000 & 20,000 \\
 & \underline{42,860} &  & \underline{42,860} \\
\hline
\end{tabular}\par\smallskip
B died on February 28\textsuperscript{th} , 2016 and according to partnership agreement his executor is entitled to be paid out as follows:\\
(a) The capital to his at the time of his death and the interest up to the time of his death at 6\% per annum.\\
(b) His appropriate share of general reserves.\\
(c) His share of profit to the date of his death which is to be taken on the basis of proceeding years’ profits.\\
(d) His share of goodwill which is calculated at two year’s purchase of the average profit of the preceding three years.\\
The investments were sold for Rs. 16,020 and B’s executor was paid off. The profits in the three years were 2013 – Rs. 7800, 2014 - Rs.9000, 2015 Rs. 9600.\\
Pass the journal entries and write the account of B.\par
}
\lqn{106}{A and B are partners sharing profits and losses as 3:2 respectively. On 31\textsuperscript{st} December 2016 they decided to dissolve the firm. The position of the firm on that date was as follows:\par
\par\smallskip\noindent\hspace*{1.2em}\begin{tabular}{|llll|}
\hline
Assets & Rs. & Liabilities & Rs. \\
Sundry assets & 50,000 & Creditors & 10,000 \\
 &  & A’s capital   25,000 &  \\
 &  & B’s capital   15,000 & 40,000 \\
 & \underline{50,000} &  & \underline{50,000} \\
\hline
\end{tabular}\par\smallskip
Assets were sold for Rs. 60,000 and creditors were paid in full. Expenses on realisation were Rs. 4000.\par
}
\lqn{107}{he following was the balance sheet of Rashid and Akbar as on the 31st December 2016.\par
\par\smallskip\noindent\hspace*{1.2em}\begin{tabular}{|lrlr|}
\hline
Assets & Rs. & Liabilities & Rs.\\
Cash at bank & 19,000 & Sundry creditors & 38,000\\
Stock-in-trade & 44,000 & Mrs. Rashid\textquoteright s loan & 10,000\\
Sundry debtors \ 20,000 & & Akbar\textquoteright s loan & 20,000\\
Less: provisions 1,000 & 19,000 & Rashid\textquoteright s capital & 10,000\\
Furniture and fixture & 4,000 & Akbar\textquoteright s capital & 8,000\\
 & \underline{86,000} & & \underline{86,000}\\
\hline
\end{tabular}\par\smallskip
The firm was dissolved on 31\textsuperscript{st}  December 2016 and the following was the result:\\
(a) Rashid agreed to pay loan of his wife.\\
(b) The assets realised the following:\par
\par\smallskip\noindent\hspace*{1.2em}\begin{tabular}{|lcl|}
\hline
 & Rs. &  \\
Stock & 37,000 &  \\
Debtors & 18,000 & The expenses of dissolution \\
Furniture and fixture & 4,500 & were Rs. 1,100 \\
\hline
\end{tabular}\par\smallskip
(c) The sundry creditors were paid off less 2\textsuperscript{1/2}  \% discount. Rashid and Akbar shared profits and losses in the ratio of 3:2. Journalize the entries to be made on the dissolution and show the realisation accounts, cash accounts and partners’ capital accounts.\par
}
\lqn{108}{The balance sheet of the firm showed the following position on 31\textsuperscript{st}  December 2016.\par
\par\smallskip\noindent\hspace*{1.2em}\begin{tabular}{|lrlr|}
\hline
Assets & Rs. & Liabilities & Rs.\\
Buildings & 40,000 & Partner capital &\\
Investments & 10,000 & Accounts: &\\
Debtors & 5,000 & D \ 32,500 &\\
Cash balance & 15,000 & E \ 27,500 & 60,000\\
 & \textbf{\underline{70,000}} & & \textbf{\underline{70,000}}\\
\hline
\end{tabular}\par\smallskip
The partnership was dissolved on 31\textsuperscript{st} December 2016. Creditors were paid at 5\% discount. D agreed to take over the building at Rs. 45,000, E took over the investments at Rs, 26,000and debtors at Rs. 3000.\\
Prepare realizing account and close the books of the firm.\par
}
\lqn{109}{A and B carrying a business as merchants and sharing profits and losses in the ratio of 1:1, dissolved their firm as on December 2016 on which date their balance sheet was as follows:\par
\par\smallskip\noindent\hspace*{1.2em}\begin{tabular}{|llll|}
\hline
Assets & Rs. & Liabilities & Rs. \\
Cash balance & 20,000 & Sundry creditors & 60,900 \\
Stock & 70,000 & Reserve funds & 21,000 \\
Premises & 90,900 & Capital accounts: &  \\
 &  & A               49,500 &  \\
 &  & B               49,500 & 99,000 \\
 & \underline{180,900} &  & \underline{180,900} \\
\hline
\end{tabular}\par\smallskip
The assets except cash at bank were sold to a company which paid Rs. 1,95,000in cash. Sundry creditors were paid Rs, 58,500 in full settlement. The expenses amounted Rs. 9,000.\\
Prepare the realisation account, cash account and partners’ accounts.\par
}
\lqn{110}{The following was the balance sheet of Waqar and Rehman as on 31\textsuperscript{st} December 2016. On that date they want to dissolve the firm.\par
\par\smallskip\noindent\hspace*{1.2em}\begin{tabular}{|lclc|}
\hline
Assets & Rs. & Liabilities & Rs. \\
Furniture & 400 & Reserve & 2000 \\
Stock & 7200 & Sundry creditors & 3800 \\
Debtors & 12000 & Capitals: &  \\
Machinery & 1800 & Waqar            11,000 &  \\
Profit \& loss a/c & 3400 & Rehman           8,000 & 19,000 \\
 & \underline{24,800} &  & \underline{24,800} \\
\hline
\end{tabular}\par\smallskip
Debtors released with a 7\textsuperscript{1/2} 5 expenditures. Stock, machinery and furniture were sold for Rs. 6000, Rs. 1400 and Rs. 500 respectively. Basing on the above balance sheet prepare the necessary ledger accounts. Partners share the profits equally.\par
}
\lqn{111}{Mr. Saeed and Mr. Ijaz are partners sharing profits and losses as 2:1. The position of the firm on 31\textsuperscript{st} December 2016 when they decided to dissolved the firm was as follows:\par
\par\smallskip\noindent\hspace*{1.2em}\begin{tabular}{|lcll|}
\hline
Assets & Rs. & Liabilities & Rs. \\
Plant \& machinery & 25,000 & Sundry creditors & 15,000 \\
Furniture & 4000 & General reserve & 10,000 \\
Stock & 10,000 & Capital accounts: &  \\
Sundry debtors & 20,000 & Mr. Saeed      22,000 &  \\
Cash at bank & 10,000 & Mr. Ijaz            22,000 & 44,000 \\
 & \underline{69,000} &  & \underline{69,000} \\
\hline
\end{tabular}\par\smallskip
The realisation shows the following results:\\
(a) Mr. Saeed took over plant and machinery and furniture at book value less 10\% per cent.\\
(b) Mr. Ijaz took over the stock and goodwill at Rs. 17,500.\\
(c) Sundry debtors realised Rs. 18,500\\
(d) Sundry creditors were settled at a discount of 5\% per cant.\\
Close the book of the firm.\par
}
\lqn{112}{A, B and C carrying on business and sharing profits and losses in ratio of 3:2:1 respectively agreed to dissolve their partnership firm on 31\textsuperscript{st} December 2016. The balance sheet of the firm as on that date is as follows:\par
\par\smallskip\noindent\hspace*{1.2em}\begin{tabular}{|llll|}
\hline
Assets & Rs. & Liabilities & Rs. \\
Machinery & 100,000 & Creditors & 88,000 \\
Stock & 60,000 & Capital accounts: &  \\
Debtors & 78,000 & A & 100,000 \\
Cash at bank & 32,000 & B & 10,000 \\
 &  & C & 72,000 \\
 & \underline{2,70,000} &  & \underline{2,70,000} \\
\hline
\end{tabular}\par\smallskip
A agreed to take the machinery at on agreed value of 80,000. A sum of Rs. 63,000 could be realised from debtors and stock was sold for Rs. 46,000.\\
The creditors were satisfied by payment of Rs. 84,000. A agreed to bear all expenses of dissolution. For this service A is to be paid  Rs. 3000. Actual expenses amounting Rs. 5000.\\
Draw up a realisation account, bank account and partners’ capital account  assuming that all partners all dissolved.\par
}
\lqn{113}{Sun and Moon sharing the profits and losses as the ratio of 4:3 decide to wind up the partnership business on 31\textsuperscript{st}  December 2016 on which date sun’s capital was 12,000 Moon’s capital was Rs. 300. As per the book there was sundry creditors outstanding to the extent of Rs. 3000 and there was a cash balance of Rs. 300. Expenses of winding up came to Rs. 250. Sundry assets realised Rs. 13,150.\\
Close the book of the firm.\par
}
\lqn{114}{X, Y and Z are partners sharing profits and losses in the proportion of 5:3:2. they decided to dissolved the partnership from on 1.,1,2016 when the balance sheet noticed as below:\par
\par\smallskip\noindent\hspace*{1.2em}\begin{tabular}{llll}
Assets & Rs. & Liabilities & Rs. \\
Cash in hand & 18000 & Sundry creditors & 1,82,000 \\
Sundry debtors & 1,52,000 & Capital accounts: &  \\
Stock-in-trade & 60,000 & X                    3,30,000 &  \\
Furniture \& fixture & 10,000 & Y                    1,30,000 & 4,60,000 \\
Plant \& machinery & 2,90,000 &  &  \\
Land \& buildings & 80,000 &  &  \\
Z’s capital account & 32,000 &  &  \\
 & \underline{\underline{6,42,000}} &  & \underline{\underline{6,42,000}} \\
\end{tabular}\par\smallskip
The assets realised: land \& building Rs. 100,000, plant \& machinery Rs. 2,50,000, sundry debtors Rs. 1,21,000, stock Rs. 47,000; furniture \& fixture Rs. 9200. Expenses of dissolution Rs. 3600.\\
Draw up the realisation account and partners’ Capital accounts.\par
}
\lqn{115}{The balance sheet of A, B and C showedas under, when they decided to dissolved their business.\par
\par\smallskip\noindent\hspace*{1.2em}\begin{tabular}{|lccc|}
\hline
Assets & Rs. & Liabilities & Rs. \\
Cash at bank & 500 & Creditors & 7000 \\
Debtors & 4500 & A’s Capital & 12000 \\
Stock & 16000 & B’s Capital & 9000 \\
Machinery & 10000 & C’s Capital & 3000 \\
 & \underline{31,000} &  & \underline{31,000} \\
\hline
\end{tabular}\par\smallskip
Debtors, stock and machinery realise 50\% of books value and expenses of realisation come to Rs.\\
50. C is insolvent and is unable to bring in anything in respect of his debt to the firm. Show that final adjustments of the account, (a) ignoring the decision in Garner v. Murray, and (b) in accordance with the decision of Garner v. Murray.\par
}
\lqn{116}{The position of Zahid, Shahid and Khalil on 30\textsuperscript{th}  June 2016 was as follows:\par
\par\smallskip\noindent\hspace*{1.2em}\begin{tabular}{llll}
Assets & Rs. & Liabilities & Rs. \\
Cash & 10,000 & Sundry Creditors & 25,200 \\
Sundry Assests & 68,000 & Zahid Load Account & 16,000 \\
Khalil Capital & 31,200 & Zahid Capital & 25,600 \\
 &  & Shahid Capital & 14,400 \\
 &  & Profit and loss a/c & 28,000 \\
 & \underline{\underline{1,09,200}} &  & \underline{\underline{1,09,200}} \\
\end{tabular}\par\smallskip
Profits and losses shared – Zahid 18/35; Shahid 7/35; Khalid 10/35. The firm is dissolved on the above date. Sundry assets realise Rs. 56,000. Sundry creditors are paid Rs. 24,000 in full settlement. Expenses amount to Rs. 3200 Khalid is insolvent.\\
Assume the capitals are not fixed. Close the books of the firm.\par
}
\lqn{117}{A, B and C are partners sharing profits and losses in the proportion of one – half, one – third and one – sixth respectively. The firm balance sheet as on 31\textsuperscript{st} December 2016 was as follows:\par
\par\smallskip\noindent\hspace*{1.2em}\begin{tabular}{|lclc|}
\hline
Assets & Rs. & Liabilities & Rs. \\
Cash at bank & 14,000 & Creditors & 31,500 \\
Debtors & 32,600 & Loan on mortgage & 4,500 \\
Stock & 22,400 & A’s Capital & 30,000 \\
Fixed assets & 15,000 & B’s Capital & 20,000 \\
C’s Capital overdrawn & 2,000 &  &  \\
 & \underline{86,000} &  & \underline{86,000} \\
\hline
\end{tabular}\par\smallskip
It was resolved the partnership as on this date. The fixed assets and stock realised Rs. 35,000. Bad debts and discounts allowed amounted Rs. 1,600. The mortgage loan was paid off. The creditors paid less discounts of 5\%.\\
There was a claim for damages pending against the firm which was settled Rs. 5,600. The expenses of realisation was amounted to Rs. 375.\\
After realisation of assets C was adjudicated an as insolvent and a dividend of 25 paise in rupee was received in respect of the firm’s claim against his estate.\\
You are required to close the books of the firm in accordance with the decision in Garner vs. Murray.\par
}
\lqn{118}{A, B and C are partners in a firm. They shared  profits and losses: A 40\% B 30\% and C 30\%. The firm was dissolved and B was appointed to realise the assets and distribute the proceeds. B is to receive 5\% commission on the amounts realised from sale of assets and to bear all expenses or realisation.\\
The balance sheet on the date of dissolution was as under:\par
\par\smallskip\noindent\hspace*{1.2em}\begin{tabular}{lccc}
Assets & Rs. & Liabilities & Rs. \\
Cash & 450 & Creditors & 17,70 \\
Debtors                 13,650 &  & A’s capital & 9,000 \\
Less: Provision         750 & 12,900 & B’s capital & 6,000 \\
Stock & 18,000 &  &  \\
C’s capital overdrawn & 1,350 &  &  \\
 & \underline{\underline{32,700}} &  & \underline{\underline{32,70}} \\
\end{tabular}\par\smallskip
Debtors realised Rs. 10,500, stock Rs. 13,500, goodwill Rs. 600. Creditors were paid Rs. 17,250 in full settlement. In addition, outstanding creditors, Rs. 150, were also paid, the expenses amounted to Rs. 180. A and B agreed to receive Rs. 900 in full settlement from C. Show the realisation account, cash accounts, and capital accounts of partners.\\
NOTE:  When one partner becomes insolvent, Garner Vs. Murray case to be applied. i.e. the debit balance in insolvent partner’s capital Account is to be debited to solvent partner’s capital accounts in their final capital ratio. However, in order to calculate final capital ratio, the abnormal items like commission on realisation of assets, being expenses etc. are to be ignored.\par
}
\lqn{119}{The balance sheet of O, P, Q and R showed the following position on dissolution:\par
\par\smallskip\noindent\hspace*{1.2em}\begin{tabular}{|lcll|}
\hline
Dr. &  & Cr. &  \\
Assets & Rs. & Liabilities & Rs. \\
Cash at bank & 34,000 & Creditors & 10,000 \\
Q’s Capital a/c & 10,000 & O’s capital & 15,000 \\
R’s Capital a/c & 3000 & P’s Capital & 10,000 \\
 &  & Realisation profit & 12,000 \\
 & \underline{47,000} &  & \underline{47,000} \\
\hline
\end{tabular}\par\smallskip
Show the final adjustment among the partners assuming that R is insolvent\par
}
\lqn{120}{X, Y and Z were partners. Their balance sheet stood as under on the date when the firm was dissolved.\par
\par\smallskip\noindent\hspace*{1.2em}\begin{tabular}{llll}
Assets & Rs. & Liabilities & Rs. \\
Sundry assets & 55,000 & Sundry creditors & 60,000 \\
Profits \& losses accounts & 12,000 & X’s capital a/c & 22,000 \\
Y’s capital a/c & 25,000 & Z’s capital a/c & 10,000 \\
 & \underline{\underline{92,000}} &  & \underline{\underline{92,000}} \\
\end{tabular}\par\smallskip
The assets realised Rs. 40,000. The expenses realisation amounted Rs. 1000. The position of the partners was as under:\par
\par\smallskip\noindent\hspace*{1.2em}\begin{tabular}{|cll|}
\hline
 & Private estates & Private liabilities \\
X & 18,000 & 20,000 \\
Y & 12,000 & 21,000 \\
Z & 12,000 & 10,000 \\
\hline
\end{tabular}\par\smallskip
Prepare the necessary ledger accounts to close the books of the firm.\par
}
\lqn{121}{Furqan and Asghar are in equal partnership. Their balance sheet stood as under:\par
\par\smallskip\noindent\hspace*{1.2em}\begin{tabular}{|lclc|}
\hline
Assets & Rs. & Liabilities & Rs. \\
Plant and machinery & 13,750 & Furqan’s capital & 6,000 \\
Furniture & 5,000 & Creditors & 39,000 \\
Debtors & 5,000 &  &  \\
Stock & 6,250 &  &  \\
Cash at bank & 3000 &  &  \\
Asghar’s drawings & 12,000 &  &  \\
 & \underline{45,000} &  & \underline{45,000} \\
\hline
\end{tabular}\par\smallskip
The partnership was dissolved and the assets were realised as follow:\\
Stock Rs. 3,500; furniture Rs. 3,000; debtors Rs. 5,000 and machinery Rs. 6,000.\\
The cost of collecting and distribution the estate amounted to Rs. 1,500. Furqan’s private estate was not sufficient even to pay his private debts, whereas in Asghar’s private estate there is a surplus of Rs. 500.\\
Prepare a realisation account, Cash account and Profit and Loss account/ Deficiency account and creditors account showing that dividend is paid to creditors.\par
}

\end{document}
